\section{The main question: which positivity can actually pass?}
\label{sec:positivity}

\subsection{Six distinct assertions}
Throughout, representations are finite dimensional over $\mathbb C$ unless
specified otherwise. Write
\[
 g(\alpha,\beta,\gamma)
 =\dim([\alpha]\otimes[\beta]\otimes[\gamma])^{\mathfrak S_d},
 \qquad
 p_{\lambda,\mu}^{\nu}
 =\dim\operatorname{Hom}_{GL(V)}
     (S_\nu V,S_\lambda(S_\mu V)).
\]
Partitions in the first formula have size $d$; in the second,
$|\nu|=|\lambda||\mu|$ and $V$ is large enough. Neither nonnegativity is
conjectural: these numbers are dimensions.

\begin{center}\small
\begin{tabularx}{\linewidth}{@{}p{.18\linewidth}X@{}}
\toprule
Assertion & Exact additional content\\
\midrule
Numerical positivity & A specified real number, determinant or norm is nonnegative.\\
Schur positivity & Expansion in the irreducible-character basis has
nonnegative integer coefficients.\\
Positive counting rule & The coefficient counts explicit combinatorial
objects; a complexity-theoretic $\#P$ rule also bounds witness length and
verification time in the stated input encoding.\\
Homogeneous PH1 & A stretching function is exactly
$\#(sP\cap\mathbb Z^D)$ for a rational polytope $P$, with the size and
separation-oracle bounds in GCT VI, Hypothesis 1.6.4.\\
Shifted counting model & The count is $\#\{z\in\mathbb Z^D:Az\le sb+c\}$.
The fixed shift $c$ changes the mathematical assertion.\\
Saturation/PH2 & GCT VI, Definitions 1.2.2 and 1.2.4 and
Hypotheses 1.6.5--1.6.6 concern constituent quasi-polynomials and bounds
on their saturation or coefficient-positivity indices. They are stronger
than values being nonnegative.\\
\bottomrule
\end{tabularx}
\end{center}

The exact stretching convention in GCT VI, Section 1.6, is
$p_{s\lambda,\mu}^{s\nu}$: the inner weight $\mu$ is fixed.
In its revised version the small-index conjectures replace a disproved
zero-index conjecture. A counterexample to naive dilation saturation must
not be presented as a counterexample to every revised hypothesis.

\subsection{Known refutations that a new purity proof cannot reverse}
\paragraph{P-K1: the literal homogeneous PH1 is false.}
Kahle--Micha\l ek, \emph{Obstructions to combinatorial formulas for
plethysm}, Theorem 5 and Remark 6, give
\[
 \phi(s)=p_{(4s),(3)}^{(7s,5s)}
 =\frac{s+r(s)}3,\qquad
 (r(0),\ldots,r(5))=(3,-1,1,0,2,-2),
 \quad r(s+6)=r(s).
 \tag{P-K1}\label{eq:km}
\]
Their Question 3 explicitly identifies this as the stretching convention
of GCT VI, Hypothesis 1.6.4. The function is not an Ehrhart function:
$\phi(1)=0$ but $-\phi(-1)=1$, contrary to Ehrhart reciprocity.
The same paper, Theorem 5, excludes one-dimensional shifted models too.
These are cited results, not new results of this investigation.

\paragraph{P-K2: the strong, zero-index Kronecker conjectures are false.}
Briand--Orellana--Rosas, \emph{Reduced Kronecker coefficients and
counter-examples to Mulmuley's strong saturation conjecture SH},
Theorem 2.4, give, in particular,
\[
 g\bigl((6s,6s),(7s,5s),(6s,4s,2s)\bigr)
 =
 \begin{cases}s/2+1&s\text{ even},\\(s-1)/2&s\text{ odd}.\end{cases}
 \tag{P-K2}
\]
The odd constituent is nonzero but vanishes at $s=1$, and has a negative
constant coefficient. This refutes the old strong SH and PH2, not the
revised polynomial index bounds. The distinction is explicit in that
paper's Section 2.5 and Mulmuley's appendix.

\paragraph{P-K3: unrestricted exact positivity is computationally hard.}
Ikenmeyer--Mulmuley--Walter, \emph{On vanishing of Kronecker coefficients},
Theorem 1.3, prove NP-hardness of Kronecker positivity, including unary
partition encoding. Fischer--Ikenmeyer, \emph{The computational complexity
of plethysm coefficients}, Theorem 3.1, prove NP-hardness for plethysm
positivity, already for fixed inner symmetric degree $3$.
Consequently, a proposed polynomial-time algorithm for all of either
family would imply $\mathsf P=\mathsf{NP}$. This is an implication;
it does not assume $\mathsf P\ne\mathsf{NP}$ and is compatible with the
supplied subexponential SAT theorem.

\subsection{A sharp repair on the counterexample ray}

\begin{proposition}[P-F1: optimal shifted model on the PH1 counterexample ray]
\label{prop:shiftedray}
For every integer $s\ge0$, the function in \eqref{eq:km} counts the integer
points of the single rational polytope
\[
 P_s=\{(t,z)\in\mathbb R^2:
 0\le z\le1,\quad 3t+4z\ge s,\quad 2t+4z\le s\}.
 \tag{P-F1}
\]
It has dimension two for $s>4$, even though $\phi$ has degree one.
Two is the smallest possible eventual dimension of an inhomogeneous
rational-polytope model for this function.
\end{proposition}
\begin{proof}
Inspection of the six residue classes gives
$\phi(s)=\lfloor s/3\rfloor+\mathbf1_{2\mid s}$, including $s=0$.
As formal power series,
\[
 \sum_{s\ge0}\phi(s)T^s
 =\frac{T^3}{(1-T)(1-T^3)}+\frac1{1-T^2}
 =\frac{1+T^4}{(1-T^2)(1-T^3)}.
\]
Expanding the last expression shows that $\phi(s)$ is the number of
integer triples
\[
 a,b\ge0,\quad z\in\{0,1\},\quad 2a+3b+4z=s.
\]
The integer coordinate change
$(a,b,z)\mapsto(k,t,z)=(2a+3b+4z,a+b,z)$ has determinant $-1$.
On $k=s$ its inverse is
$a=3t+4z-s$, $b=s-2t-4z$.
It is therefore an exact bijection of the relevant integer lattices,
and nonnegativity becomes the two displayed inequalities for $P_s$.
For $s>4$, at each $0<z<1$ the interval for $t$ has length
$(s-4z)/6>0$. The area is $(s-2)/6$, so the dimension is two.
A zero-dimensional convex polytope has at most one integer point,
whereas $\phi$ is unbounded. For completeness, the exclusion of
one-dimensional models also holds when an affine line is embedded in
a larger ambient lattice, as follows. Suppose
$Q_s=\{x\in\mathbb R^N:Ax\le sb+c\}$ is a bounded, eventually
one-dimensional model. Candidate vertices are the finitely many affine
functions of $s$ obtained by solving independent sets of $N$ active
constraints. Their feasibility and the equality of each constraint at
each vertex stabilize for sufficiently large $s$, since all these tests
are signs of affine functions. The constraints tight at every vertex
therefore give a fixed affine-hull system $Mx=sb'+c'$, whose kernel is
a rational line.

The count $\phi(s)$ is positive for every integer $s\ge2$. Choose
integer support points $x_{s_0},x_{s_0+1}$ at two consecutive sufficiently
large parameters, put $w=x_{s_0+1}-x_{s_0}$ and $v_0=x_{s_0}-s_0w$.
Then $v_0+sw$ lies in the affine support for every integer $s$.
If $v$ is a primitive integer generator of $\ker M$, all integer
points in that support are uniquely $v_0+sw+tv$, $t\in\mathbb Z$.
Substitution turns the inequalities into a bounded interval for $t$.
On a tail its endpoints are affine, say $as+\alpha$ and $bs+\beta$,
because maxima and minima of finitely many affine functions stabilize.
Its integer count differs from its length by a bounded quantity,
so $b-a=1/3$. Translate $t$ by an integer multiple of $s$ to arrange
$0\le a<1$.
If $b\le1$, at most one integer can enter at the right endpoint in a
unit step, and the nondecreasing left endpoint cannot add any.
But $\phi(s+1)-\phi(s)=2$ at arbitrarily large $s\equiv5\pmod6$.
Hence $b>1$. Then at least one integer enters at the right in every
step and at most one leaves at the left, so the count is eventually
nondecreasing. This contradicts its drop by one at every large
$s\equiv0\pmod6$. Thus no eventual dimension-one model exists.
This extends the one-variable exclusion in the cited theorem to an
arbitrary fixed ambient lattice and completes the minimality proof.
\end{proof}

This repairs this ray, not homogeneous PH1. It also gives a positive
formula for the equal rectangular Kronecker ray
$g((4s,4s,4s),(4s,4s,4s),(7s,5s))$; the equality of the two coefficient
families is recorded by Kahle--Micha\l ek on page 6.
No assertion of historical priority is made for this elementary
reconstruction. Positive formulas in the two-row range have independent
literature, including Pak--Panova--Swanson (2026 version).

\begin{proposition}[P-F2: from an explicit positive Hilbert series to one shifted polytope]
\label{prop:selector}
Let
\[
 F(T)=\frac{\sum_{i\in I}h_iT^i}{\prod_{j=1}^{r}(1-T^{a_j})},
 \qquad I\subset\mathbb Z_{\ge0}\text{ finite},\quad
 h_i,a_j\in\mathbb Z_{>0}.
\]
Every coefficient $[T^s]F(T)$, for $s\ge0$, is the integer-point count
of one bounded rational polytope whose inequalities have right sides
affine in $s$. Its description length is polynomial in the explicit
sparse numerator, the expanded list of denominator factors, and the
binary encoding of $s$.
\end{proposition}
\begin{proof}
Use real variables $x_1,\ldots,x_r,u$ and $e_i$ for $i\in I$ and impose
\[
 x_j,u,e_i\ge0,\qquad
 \sum_{i\in I}e_i=1,\qquad
 \sum_{j=1}^{r}a_jx_j+\sum_{i\in I}i e_i=s,\qquad
 u+1\le\sum_{i\in I}h_i e_i.
 \tag{P-F2}
\]
These are linear equalities and inequalities, hence define a convex
rational polytope (equalities can be replaced by two inequalities).
At an integer point, the nonnegative $e_i$ summing to one select exactly
one $i$. For that choice there are exactly $h_i$ possibilities
$u=0,\ldots,h_i-1$. The remaining equation counts nonnegative integer
solutions of $\sum_j a_jx_j=s-i$, which is the coefficient of
$T^{s-i}$ in the denominator's reciprocal. Summing the disjoint integer
selector slices gives exactly $[T^s]F$.
Boundedness holds even for real points: $0\le e_i\le1$,
$0\le u\le\max_i h_i-1$, and $0\le x_j\le s/a_j$ whenever feasible.
There are $r+|I|+1$ variables and linearly many constraints; their
coefficients are the explicitly supplied integers $a_j,i,h_i$, together
with $0,1,-1$. This proves the description bound. If the numerator is
zero, use the fixed inconsistent inequality $0\le-1$.
\end{proof}

\paragraph{What P-F2 settles.}
GCT VI, Theorem 1.6.1(d), supplies positive rational Hilbert forms for
its canonical models. P-F2 supplies an ordinary, possibly
large-dimensional shifted counting model from any such form.
The gate for an \emph{efficient uniform} model is now concrete:
construct sufficiently short positive numerator and denominator data
uniformly, or replace their explicit enumeration by an effective
description with the required separation bounds. A mere existence
theorem for an h.s.o.p.\ or a nonnegative Hilbert numerator does not give
those bounds. Neither P-F2 nor a polynomial description proves
saturation: integer feasibility in a rational polytope remains a
separate issue. This is a sufficient route to shifted counting, not a
claim that every possible route must use Hilbert series.

\subsection{Positive determinants already exist for the coefficients}

\begin{proposition}[P-F3: an unconditional determinant certificate, with its cost exposed]
\label{prop:projector}
Let $K$ be a compact group and $W$ a finite-dimensional continuous
unitary representation. With normalized Haar measure put
$P=\int_K\rho(k)\,dk$ and $m=\dim W^K$. Then
\[
 P=P^*=P^2,\qquad
 \det(I+tP)=(1+t)^m
 =\sum_{j=0}^{m}\binom mj t^j.
\]
For a finite group, the integral is the finite average.
Consequently every Kronecker coefficient and every plethysm coefficient
has a positive-semidefinite determinant representation of this form.
\end{proposition}
\begin{proof}
For every $h\in K$, invariance of Haar measure gives $\rho(h)P=P$,
so the image of $P$ lies in $W^K$. If $w\in W^K$, then $Pw=w$.
Thus $P$ is a projection onto $W^K$.
Unitarity and invariance under inversion give
$P^*=\int_K\rho(k^{-1})\,dk=P$. It has $m$ eigenvalues one
and all remaining eigenvalues zero, which proves the determinant formula.
For Kronecker take $K=\mathfrak S_d$ and
$W=[\alpha]\otimes[\beta]\otimes[\gamma]$.
For plethysm take $K=U(V)$ and
$W=(S_\nu V)^*\otimes S_\lambda(S_\mu V)$.
An algebraic map commuting with $U(V)$ commutes with $GL(V)$, because
$U(V)$ is Zariski dense in $GL(V)$, or equivalently because its complexified
Lie algebra is $\mathfrak{gl}(V)$ and $GL(V)$ is connected.
Hence $W^K$ is precisely the required Hom space.
\end{proof}

The matrix in P-F3 acts on the \emph{whole} tensor/Hom space. Its
dimension and the direct averaging cost are not polynomially bounded
in a compressed coefficient input. The determinant formula contains the
unknown rank; it is not a noncircular efficient algorithm for that rank.
It shows exactly why an additional positive determinant, without a
size bound and a coefficient-identification theorem, does not supply
PH1 or a multiplicity comparison.

\begin{proposition}[P-F4: scalar positivity does not imply character positivity]
\label{prop:scalar}
A real nonnegative class function, even an integer virtual character,
need not have nonnegative irreducible multiplicities.
\end{proposition}
\begin{proof}
On $C_2=\{1,\tau\}$ let $\mathbf1$ and $\varepsilon$ be the trivial
and sign characters. The virtual character
$f=\mathbf1-\varepsilon$ has values $f(1)=0$, $f(\tau)=2$,
both nonnegative, but its sign-character coefficient is $-1$.
Moreover the kernel $(f(x^{-1}y))_{x,y\in C_2}$ is
$\left(\begin{smallmatrix}0&2\\2&0\end{smallmatrix}\right)$,
with eigenvalues $2,-2$. Pointwise positivity is therefore different
even from positive type. The scalar inequality $W(1_S)\ge0$ cannot
be substituted for positivity of all representation projectors without
an additional identity.
\end{proof}

\begin{proposition}[P-F5: ordinary interlacing is not plethysm restriction]
\label{prop:interlacing}
The homomorphism $GL_2\to GL_3$ defined by $\operatorname{Sym}^2$
is not conjugate to the standard block embedding
$g\mapsto\operatorname{diag}(g,1)$. Their branching multiplicities
already disagree on the defining $GL_3$ representation.
\end{proposition}
\begin{proof}
The pullback of $\mathbb C^3$ by the first map is
$\operatorname{Sym}^2\mathbb C^2$, an irreducible three-dimensional
module. The pullback by the second is $\mathbb C^2\oplus\mathbb C$,
which is reducible. More explicitly, $zI_2$ acts in the first as $z^2I_3$,
and in the second with eigenvalues $z,z,1$; for example $z=2$ gives
different spectra. Conjugate homomorphisms would give isomorphic pullbacks,
so they cannot be conjugate.
\end{proof}

Thus GGP interlacing and a chain of standard unitary restrictions do not
by themselves compute restriction along the representation homomorphism
$GL(V)\to GL(S_\mu V)$. A theta-lift route would have to identify these
specific multiplicity spaces, prove compatibility of its compositions
and specializations, and control sizes. An order of vanishing gained
by an analytic lift is not, by definition, the dimension of one of
these Hom spaces. No such identifying theorem is stated in the supplied
four-page note or in the five supplied Lean declarations.

\subsection{Purity: the exact categorical gate and a complexity barrier}

\paragraph{P-K4: the established geometric template.}
Deligne's \emph{Weil II}, Theorem 3.3.1 and Corollaries 3.3.4--3.3.6,
control weights under direct image; Corollary 3.3.9 gives the expected
pure weights for smooth proper cohomology. BBD, \emph{Faisceaux pervers}, Theorems 5.3.8 and 5.4.5,
give geometric semisimplicity and decomposition for pure complexes,
and Theorem 6.2.5 gives the characteristic-zero decomposition theorem.
For Schubert varieties of a split connected reductive group over
$\mathbb F_q$, Haines, \emph{A proof of the Kazhdan--Lusztig purity
theorem via the decomposition theorem of BBD}, Theorem 2.0.1,
gives odd-stalk vanishing and the eigenvalues $q^{i/2}$ in even degree
$i$ for the unshifted intersection complex. His Theorem 5.0.4 gives
$P_{y,w}(q)=\sum_i\dim H^{2i}_{yB}(IC(X_w))q^i$, identifying the
Kazhdan--Lusztig polynomials with local intersection-cohomology
dimensions. These statements have hypotheses and identifications;
purity alone is not that identification.

\begin{proposition}[P-F6: a sufficient categorical transfer criterion]
\label{prop:categorical}
Let $\mathcal C$ be an additive semisimple graded category whose
indecomposables are $T^jL_\lambda$, with unique finite decompositions.
Suppose a bilinear product $\star$ preserves $\mathcal C$ and satisfies
$(T^iA)\star(T^jB)\simeq T^{i+j}(A\star B)$.
Suppose its split Grothendieck ring is identified with a based
$\mathbb Z[v,v^{-1}]$-algebra $A$, with $[L_\lambda]\mapsto b_\lambda$
and $[TM]\mapsto v[M]$. Then the structure constants of the
$b_\lambda$ lie in $\mathbb Z_{\ge0}[v,v^{-1}]$.
If specialization at $v=1$ identifies these constants with specified
representation multiplicities, those specializations are nonnegative
sums of actual multiplicities. This does not assert an efficient
combinatorial rule.
\end{proposition}
\begin{proof}
For each pair $\alpha,\beta$, the assumptions give an actual
decomposition
$L_\alpha\star L_\beta\simeq
\bigoplus_{\lambda,j}(T^jL_\lambda)^{\oplus m_{\alpha\beta}^{\lambda,j}}$
with $m_{\alpha\beta}^{\lambda,j}\in\mathbb Z_{\ge0}$.
Pass to the split Grothendieck group and apply the specified ring
identification. The coefficient of $b_\lambda$ is
$\sum_jm_{\alpha\beta}^{\lambda,j}v^j$.
Specialization at $v=1$ sums the nonnegative multiplicities, once the
asserted comparison with representations has independently been proved.
\end{proof}

The use of the \emph{split} Grothendieck group matters. In a triangulated
Grothendieck group a cohomological shift introduces a minus sign.
The rank-one constant Frobenius module $\mathbb Q_\ell$, with Frobenius
one, placed in cohomological degree one has alternating trace $-1$,
despite having a pure eigenvalue. Parity, basis identification, product
compatibility and geometric semisimplicity must all be supplied in the
intended realization. P-F6 does not produce an efficient description of
its multiplicity spaces.

\begin{proposition}[P-F7: even pure Tate geometry can encode all \#3SAT counts]
\label{prop:tatesat}
Fix a finite field $\mathbb F_q$ and a prime $\ell\ne\operatorname{char}\mathbb F_q$.
From any 3-CNF formula $F$ with $n$
variables and $m$ clauses one can construct, in time polynomial in
$n+m$, equations for a finite smooth proper $\mathbb F_q$-scheme $X_F$
such that
\[
 \#X_F(\mathbb F_{q^r})=\#\operatorname{SAT}(F)\quad(r\ge1),\qquad
 H^0_{\mathrm{\acute et}}(X_{F,\overline{\mathbb F}_q},\mathbb Q_\ell)
 \simeq \mathbb Q_\ell^{\,\#\operatorname{SAT}(F)}
\]
with Frobenius the identity, and all other cohomology zero.
On this explicitly presented family, nonemptiness is NP-complete and
the count is \#P-complete even under this promised purity, parity and
smoothness.
\end{proposition}
\begin{proof}
Impose $x_i^2-x_i=0$ for $1\le i\le n$.
For a positive literal $x_i$ let its false-value factor be $1-x_i$;
for a negative literal $\neg x_i$ let it be $x_i$.
For each clause impose the product of its at most three false-value
factors equal to zero.
The algebra
$\mathbb F_q[x_1,\ldots,x_n]/(x_i^2-x_i:1\le i\le n)$
is, by repeated Chinese remainder, the product of one copy of
$\mathbb F_q$ for every Boolean assignment. This remains true in
characteristic two, since $0$ and $1$ are distinct and
$x_i(x_i-1)$ has distinct roots.
In a factor corresponding to an assignment, each clause product is
zero exactly when the clause is satisfied, and otherwise equals one.
Quotienting by all clause products therefore deletes precisely the
unsatisfying factors. The resulting coordinate algebra is
$\prod_{\sigma\models F}\mathbb F_q$ (the zero algebra for no assignments).
Consequently $X_F$ is a disjoint union of rational points, is finite
\'etale, and hence smooth and proper. The displayed point counts and
cohomology follow point by point.
There are $n+m$ equations of degree at most three with coefficients of
bounded size. The construction is parsimonious, so the
NP-completeness and \#P-completeness of 3SAT transfer (a constructive
verifier proof is given in C-COUNT-03 below). Membership follows
by guessing and checking an assignment.
\end{proof}

This is a reduction, not a proof that $\mathsf P\ne\mathsf{NP}$.
It establishes an NP-hardness boundary for any proposed recipe using
a polynomial list of equations and pure even Tate cohomology:
a polynomial algorithm from those data would solve 3SAT.
The supplied ETH theorem improves the subexponential computation of
these examples, without changing the exact polynomial-time implication.

\subsection{Nonstandard quantum groups: what remains to be constructed}
GCT VII, Theorem 1.2, constructs the nonstandard quantization.
Its Conjectures 1.3 and 1.4 are duality and reciprocity assertions.
GCT VIII's Conjectures 2.6--2.9 and 3.4--3.7 impose cell, positivity,
saturation and compatibility properties. The general setting even
permits signed factors and a relaxed positivity condition; it is not
the assertion that every raw coefficient is nonnegative.
The intended geometric realization and specialization remain additional
input. P-F6 supplies a sufficient transfer criterion once those data
exist; P-F3--P-F7 identify why determinant signs, interlacing or purity
alone do not supply them. No theorem of nonexistence of such a future
realization is claimed.

\subsection{Occurrence, multiplicity and actual equations}

\paragraph{P-K5: the no-occurrence theorem has a precise range.}
B\"urgisser--Ikenmeyer--Panova, \emph{No occurrence obstructions in
geometric complexity theory}, Theorem 1.4, prove that for
$n\ge m^{25}$ every irreducible occurring in the coordinate ring of
the padded permanent orbit closure also occurs in the determinant
orbit closure. This is a theorem about these closures, not merely an
upper bound using rectangular Kronecker coefficients. It does not
exclude multiplicity obstructions or all other uses of representation
theory in complexity.

\begin{proposition}[P-F8: the surviving multiplicity test and its logical limit]
\label{prop:multiplicity}
Let $X,Y$ be $G$-stable affine cones in one representation of a complex
reductive group, and let $m_X(d,\lambda)$, $m_Y(d,\lambda)$ be their
graded coordinate-ring multiplicities. If
$m_Y(d,\lambda)>m_X(d,\lambda)$ for some $(d,\lambda)$ then
$Y\not\subseteq X$. The reverse inference from all weak inequalities
to containment is false, even when all the multiplicities agree.
In the no-occurrence range above, any multiplicity obstruction for the
specified determinant and permanent closures has both multiplicities
strictly positive.
\end{proposition}
\begin{proof}
If $Y\subseteq X$, restriction is a surjective graded $G$-map
$\mathbb C[X]\twoheadrightarrow\mathbb C[Y]$.
Finite-dimensional rational $G$-modules in characteristic zero are
semisimple, so taking the $\lambda$-multiplicity in degree $d$ gives
$m_Y(d,\lambda)\le m_X(d,\lambda)$. This proves the first assertion.
For the failure of the converse, let $G=\mathbb G_m$ act on
$\mathbb A^2$ by scalar multiplication, let $X$ be the first coordinate
axis and $Y$ the second. Both are stable irreducible cones and both
coordinate rings have exactly one copy of the character of weight
$-d$ in each degree $d$. Nevertheless neither axis contains the other.
Finally a strict inequality implies $m_Y>0$; P-K5 then gives $m_X>0$.
\end{proof}

The theorem sought for determinant versus permanent is consequently
a strict comparison of the \emph{actual quotient} multiplicities, or
an explicit equation separating the closures. Ambient plethysm
multiplicities and determinant-orbit Kronecker bounds do not determine
those quotient multiplicities. The Chow/Foulkes--Howe map in the
plethysm section makes the missing rank information explicit.

\subsection{Disposition of the main request}
The literal homogeneous PH1 and the old strong zero-index assertions
are refuted by the cited counterexamples. A shifted two-dimensional
model for the PH1 counterexample is proved in P-F1; the general
positive-rational-form transfer is P-F2. P-F11--P-F12 prove the
quantitative implication from revised PH3 to revised PH2, and under
PH3 bound the size of a shifted model in the expanded group parameters.
The weaker positive-form clause and uniform construction remain the
gate on that route. P-F13 makes reconstruction and positive-form verification
finite under a stated denominator bound. P-F14--P-F15 refute the
printed rank-only PH3, PH2 and saturation-index bounds for general
subgroup restriction. P-F16--P-F18 instead prove finite bounds for
fixed semisimple embeddings and fixed inner dimension in ordinary
plethysm. P-F19 proves that the optimal canonical envelope has an
exponential subsequence, refuting the stronger canonical clause of
ordinary plethysm PH3. It leaves the numerical positive-form clause,
PH2 and the narrower saturation bounds unresolved.
PLE-F7--PLE-F8 identify the actual canonical multiplicity algebra on
the obstructing ray as a normal surface with pure even Tate ordinary
and intersection cohomology and a positive multiplicative basis.
Its graded Hilbert function still violates homogeneous PH1.
The known geometric route to
canonical positivity can be expressed by P-F6, but its nonstandard
realization, efficiency bounds and comparison with the relevant
coordinate-ring quotients have not been constructed here.
The supplied RH/Dirichlet-GRH statements do not state any of those
constructions. Reproving an existing finite-field purity theorem does
not enlarge its domain of applicability to an unidentified object.
This is a precise account of a missing deduction, not a logical
independence theorem between RH and a GCT conjecture.
